It is now easy to give applications of the implicit function theorem to elliptic PDEs. Consider the equation
\begin{equation}\label{eq:semilinear-small-data-dirichlet-equation}
\Delta u+f(u)=g(\bm{x})
\end{equation}
on a bounded smooth domain in $\mathbb{R}^3$, subject to Dirichlet boundary conditions. Assume that $f$ is of class $C^3$ with $f(0)=f^{\prime}(0)=0$. Then it follows from Lemma~\ref{lem:sobolev-pointwise-composition-differentiability} that the mapping $u\mapsto f(u)$ is continuously differentiable from $H^2(\Omega)$ into itself. We can now apply the inverse function theorem to conclude that, for sufficiently small $g\in L^2(\Omega)$, the equation has a solution $u\in H^2(\Omega)\cap H_0^1(\Omega)$. Moreover, this solution is unique among solutions of small norm. (We note that actually we only need the mapping $u\mapsto f(u)$ to be continuously differentiable as a mapping from $H^2(\Omega)$ to $L^2(\Omega)$. For this, it would suffice that $f$ be of class $C^1$.) It is clear how to formulate and prove more general results involving nonlinearities which depend on derivatives of $u$ and/or nonlinear boundary conditions.

As an example of a bifurcation problem, consider the equation
\begin{equation}\label{eq:semilinear-dirichlet-eigenvalue-bifurcation-equation}
\Delta u+\lambda u=f(u),
\end{equation}
again with Dirichlet boundary condition and with $\Omega$ and $f$ as before. The zero solution is the only small solution as long as $\lambda$ is not an eigenvalue of $-\Delta$. If $\lambda_0$ is a simple eigenvalue, the assumptions of the previous section apply, since the range of $\Delta+\lambda$ is precisely the orthogonal complement of the nullspace, and, with $u_0$ denoting the eigenfunction, and $F(u,\lambda)=\Delta u+\lambda u-f(u)$, we find $D_{u\lambda}F(0,\lambda_0)u_0=u_0$. It turns out that the first eigenvalue of the Dirichlet problem for the Laplacian is always simple.

\begin{Lemma}
\label{lem:first-dirichlet-eigenvalue-simplicity}
Let $\Omega$ be a bounded domain. Then the first eigenvalue of the Dirichlet problem for Laplace\textquotesingle s equation is simple.
\end{Lemma}

\begingroup
\renewcommand{\endProof}{\par\nobreak\hfill$\blacksquare$\par\vspace{0.1cm}}
\begin{Proof}
Let $\lambda_0$ be the smallest eigenvalue; then we have
\begin{equation}\label{eq:first-dirichlet-eigenvalue-printed-norm-minimum}
\lambda_0=\min_{\substack{u\in H_0^1(\Omega),\|u\|_2=1}}\|u\|_{1,2}^2.
\end{equation}
Recall Problem~\ref{ex:sturm-liouville-smallest-eigenvalue-variational-characterization}, where the analogous result was established for Sturm-Liouville problems. Let $u_0$ be a minimizer, and assume that $u_0$ has zeros inside $\Omega$. We can now construct a sequence of piecewise polynomials $p_n$, which converges to $u_0$ in $H_0^1(\Omega)$, and uniformly on compact subsets of $\Omega$. Clearly, we have
\begin{equation}\label{eq:first-dirichlet-eigenvalue-piecewise-polynomial-quotients}
\lambda_0=\lim_{n\to\infty}\|p_n\|_{1,2}^2/\|p_n\|_2^2.
\end{equation}
Now $|p_n|$ has the same $H^1$- and $L^2$-norms as $p_n$ itself; in particular, the $H^1$-norms of $|p_n|$ are bounded as $n\to\infty$. Hence there is a subsequence of $|p_n|$ which converges weakly in $H_0^1(\Omega)$, and strongly in $L^2(\Omega)$. Since the $p_n$ converge to $u_0$ uniformly on compact subsets of $\Omega$, it also follows that $|p_n|\to|u_0|$. We thus find that
\begin{equation}\label{eq:first-dirichlet-eigenfunction-absolute-value-quotient}
\bigl\||u_0|\bigr\|_{1,2}^2/\bigl\||u_0|\bigr\|_2^2\leq\lambda_0,
\end{equation}
i.e., $|u_0|$ must also be an eigenfunction. But $|u_0|$ is non-negative and has zeros inside $\Omega$, which contradicts the maximum principle.

Hence our assumption of an eigenfunction that changes sign must have been wrong, and every eigenfunction is either positive or negative. No two such functions can be orthogonal to each other, and hence the eigenvalue is simple.
\end{Proof}
\endgroup
